\documentclass{article}
\usepackage{amsmath, amssymb, amsthm}

\newtheorem{theorem}{Theorem}[section]
\newtheorem{lemma}[theorem]{Lemma}
\theoremstyle{definition}
\newtheorem{definition}[theorem]{Definition}

\title{Infinitude of Primes}
\author{Fermat Demo}
\date{}

\begin{document}
\maketitle

% ── Open this file in Fermat and click "Submit All" to prove all three results.

\section{Foundations}

\begin{definition}[Prime Number]
\label{def:prime}
A natural number $p > 1$ is \emph{prime} if its only positive divisors are
$1$ and $p$ itself.
\end{definition}

\begin{definition}[Composite Number]
\label{def:composite}
A natural number $n > 1$ is \emph{composite} if it is not prime, i.e.\ there
exist integers $a, b$ with $1 < a, b < n$ such that $n = ab$.
\end{definition}

\begin{lemma}[Every integer $> 1$ has a prime divisor]
\label{lem:prime-divisor}
For every integer $n > 1$ there exists a prime $p$ such that $p \mid n$.
\end{lemma}
% [PROVE IT: Easy]

\section{Main Result}

\begin{theorem}[Infinitude of Primes]
\label{thm:inf-primes}
There are infinitely many prime numbers.
\end{theorem}
% [PROVE IT: Easy]
% SKETCH: Assume finitely many primes p_1,...,p_k. Consider N = p_1*...*p_k + 1.
%   By Lemma~\ref{lem:prime-divisor}, N has a prime divisor p.
%   But p cannot be any p_i (it would divide N - p_1...p_k = 1). Contradiction.

\begin{lemma}[Euclid's Lemma]
\label{lem:euclid}
If a prime $p$ divides a product $ab$ of integers, then $p \mid a$ or $p \mid b$.
\end{lemma}
% [PROVE IT: Medium]

\end{document}
